\originalchapter{综合练习与逐讲对应}
\sourcelecture{15,28,37--39}
前两次复习讲义分别把第1--14讲与第17--27讲的内容重新串联. 其中的计数、公理、分布与矩计算已在前文推导, 不把同一证明按课件渐显页重复排印. 最后三次复习强调多个工具的配合: 先用指标拆计数, 再检查条件分布, 或把停止问题、极限定理和定价计算联系起来. 本章按这些训练目标设置完整练习, 参数与情境另行组织.

\originalsection{指标计数与条件分布}
\begin{exercise}
两次独立排名均在 $n!$ 个排列中均匀选择, $n\ge2$. 令 $F$ 为名次不变的对象数, $U_s$ 为恰好前进 $s$ 位的对象数, $1\le s<n$. 求二者期望与方差. 后退 $s$ 位的对象数又如何?
\end{exercise}
\begin{solution}
固定第一次排名, 第二次仍为均匀排列. 令 $I_i$ 表示原第 $i$ 名最后仍第 $i$ 名. 有 $\E I_i=1/n$, 对 $i\ne j$, $\E I_iI_j=1/[n(n-1)]$. 因而
\[
\E F=1,\qquad \E F^2=n\frac1n+n(n-1)\frac1{n(n-1)}=2,
\qquad\Var F=1.
\]
对前进数, 只有原名次 $s+1,\ldots,n$ 可以前进 $s$ 位. 记 $m=n-s$, 相应指标为 $J_i=\ind_{\{\text{新名次}=i-s\}}$. 任意两个要求的目的名次不同, 所以单项与双项的概率仍是 $1/n$ 与 $1/[n(n-1)]$. 得
\[
\E U_s=\frac mn,\qquad
\Var U_s=\frac mn+\frac{m(m-1)}{n(n-1)}-\frac{m^2}{n^2}.
\]
后退情形同理, 可行原名次为 $1,\ldots,n-s$, 因而期望与方差相同. 例如 $n=10,s=3$ 时, $\E U_3=7/10$, $\Var U_3=203/300$. 指标之间并不独立; 直接代入二项方差会丢掉排列的约束.
\end{solution}

\begin{exercise}
十二次独立抽签, 每次等概率取得五种符号之一. 已知第三种符号恰好出现三次, 求第一种符号恰好出现五次的条件概率.
\end{exercise}
\begin{solution}
先固定第三种符号所在的三个位置. 条件下其余九个位置独立且均匀地取其余四种符号, 所以第一种符号的计数为 $\Bin(9,1/4)$, 答案是
\[
\binom95\left(\frac14\right)^5\left(\frac34\right)^4.
\]
也可直接用比值核验. 交集概率是 $\binom{12}{3}\binom95 3^4/5^{12}$, 条件事件概率是 $\binom{12}{3}4^9/5^{12}$, 相除得到同一值. 不指定三个位置时仍可求这个二项分布, 因为每一种位置安排给出同样的条件计数分布.
\end{solution}

\originalsection{连续分布与独立和}
\begin{exercise}
两类报警分别是速率为2与5的独立泊松过程, 时间单位为一个月. 令 $A,B$ 是各自首次报警时间. 求 $\E A^2$, $\Cov(A,B)$, $\min(A,B)$ 的密度, 以及一个月分成六个等长区间后, 没有任何报警的区间数的期望和方差.
\end{exercise}
\begin{solution}
等待时间分别为 $\Exp(2)$ 与 $\Exp(5)$, 独立性给 $\Cov(A,B)=0$. 指数矩公式给 $\E A^2=2!/2^2=1/2$. 对 $t\ge0$,
\[
\Prob(\min(A,B)>t)=\Prob(A>t)\Prob(B>t)=\ee^{-7t},
\]
所以最小值密度为 $7\ee^{-7t}$, $t>0$, 其他处为0. 每个区间无报警的概率为 $p=\ee^{-7/6}$. 不同区间的两类增量共同独立, 故无报警区间数 $K\sim\Bin(6,p)$, 从而 $\E K=6\ee^{-7/6}$, $\Var K=6p(1-p)$. 只有期望时, 无须先证明六个指标独立; 求这个方差时独立增量才真正派上用场.
\end{solution}

\begin{exercise}
$X$ 在 $[-2,2]$ 均匀分布. 求 $\Var(X^2)$. 再令 $X_1,\ldots,X_n$ 为独立同分布样本, 求最小值 $M$ 的密度和期望.
\end{exercise}
\begin{solution}
对正整数 $k$, 对称积分给 $\E X^{2k}=2^{2k}/(2k+1)$. 因而
\[
\Var(X^2)=\frac{16}{5}-\left(\frac43\right)^2=\frac{64}{45}.
\]
当 $-2<x<2$ 时, $\Prob(M>x)=((2-x)/4)^n$, 求导得
\[
f_M(x)=\frac n4\left(\frac{2-x}{4}\right)^{n-1},
\]
支持集之外为0. 为算期望, 令 $V=(M+2)/4$, 则 $V$ 是 $n$ 个 $(0,1)$ 均匀样本的最小值. 尾积分给 $\E V=\int_0^1(1-t)^n\dd t=1/(n+1)$. 所以 $\E M=-2+4/(n+1)$. 大样本的最小值向左端点集中, 和单个样本均值0是不同问题.
\end{solution}

\begin{exercise}
设 $X_1,\ldots,X_n$ 独立同分布, $X$ 的矩母函数 $M_X(t)$ 在零的某邻域有限. 求样本平均 $\overline X_n$ 的矩母函数. 解释标准柯西样本平均为什么不会依概率收敛到一个有限常数.
\end{exercise}
\begin{solution}
指数把和变成积, 独立性再使期望分解:
\[
M_{\overline X_n}(t)=\E\prod_{i=1}^n\ee^{tX_i/n}
=\bigl(M_X(t/n)\bigr)^n.
\]
标准柯西分布的特征函数为 $\ee^{-|t|}$, 因而柯西平均的特征函数仍为 $(\ee^{-|t|/n})^n=\ee^{-|t|}$, 即它仍具有同一柯西分布. 若收敛到常数 $c$, 则对任意 $\eps>0$ 都应有 $\Prob(|\overline X_n-c|>\eps)\to0$. 但这个概率不随 $n$ 改变且为正, 因为柯西密度在整条实轴为正. 不存在这样的收敛. 大数定律的可积条件不能由独立同分布替代.
\end{solution}

\originalsection{熵、状态与时间平均}
\begin{exercise}
$X,Y$ 独立, 均以概率 $p\in(0,1)$ 取1, 以概率 $1-p$ 取0. 比较 $H(X+Y)$ 与 $H(X,Y)$, 并求差值. 证明对任意有限离散变量, $H(X+Y)\le H(X,Y)$.
\end{exercise}
\begin{solution}
记 $q=1-p$, 与第11章相同, 用以2为底的对数计熵. $X+Y$ 的概率分别为 $q^2,2pq,p^2$, 而联合分布的四个概率为 $q^2,pq,pq,p^2$. 因而
\[
H(X,Y)=2[-p\log_2 p-q\log_2 q],\qquad
H(X,Y)-H(X+Y)=2pq>0.
\]
相加只在和为1时丢失``哪个变量取1''的信息, 两种可能此时等概率, 所以差值正是这个条件下的一比特信息乘发生概率. 一般设 $Z=X+Y$. 因 $Z$ 是 $(X,Y)$ 的函数, 链式公式给
\[
H(X,Y)=H(Z)+H(X,Y\mid Z)\ge H(Z).
\]
有限状态保证所有熵有限; 若扩展到可数状态, 仍可按非负求和建立不等式, 但不应把两个无穷熵形式相减.
\end{solution}

\begin{exercise}
仓库有两件可重复使用的工具. 每天有两个操作各执行一次, 次序由独立公平硬币决定. 操作 $D$ 在架上有工具时移走一件, 无工具时不改变; 操作 $R$ 在架上至少有一件时不改变, 在架上没有时把两件都放回. 以每天开始时架上数量为状态. 求转移矩阵、稳态分布, 以及长期 $D$ 遇到空架的日比例.
\end{exercise}
\begin{solution}
状态依次为0,1,2. 逐个复合两种操作, $RD$ 把状态映为 $(1,0,1)$, $DR$ 把状态映为 $(2,2,1)$. 所以行向量约定 下
\[
P=\begin{pmatrix}0&1/2&1/2\\1/2&0&1/2\\0&1&0\end{pmatrix}.
\]
令 $\pi P=\pi$, 解得 $\pi_0=\pi_1/2$, $\pi_2=(\pi_0+\pi_1)/2$, 再用和为1得 $\pi=(2/9,4/9,1/3)$. 所有状态互达, 状态0有长度2和3的回路, 其公约数为1, 故有限状态不可约非周期收敛定理适用.

$D$ 遇空架只会在日初状态为0且 $D$ 先执行时发生; 若 $R$ 先执行就已有两件. 这个日事件的稳态概率为 $\pi_0/2=1/9$. 为得到实际时间比例, 将日初数量 $X_n$ 与当天操作次序 $C_n$ 一起作为六状态链. 从 $(x,c)$ 转到 $(g_c(x),d)$ 的概率为 $1/2$, 其中 $d$ 是次日的新次序. 原来的数量链不可约, 未来各次序都可正概率指定, 所以扩大后的链也不可约. 分布 $\widetilde\pi(x,c)=\pi_x/2$ 平稳: 对任意 $(y,d)$, 流入质量为 $\tfrac12\sum_x\pi_xP_{xy}=\pi_y/2$. 对这个六状态链的奖励函数 $f(x,c)=\ind_{\{x=0,\ c=D\text{先}\}}$ 使用第10章的时间平均定理, 得实际长期日比例几乎必然为 $\sum_{x,c}\widetilde\pi(x,c)f(x,c)=1/9$. 因而这里得到的是路径上的比例, 还不只是某个遥远日的事件概率.
\end{solution}

\originalsection{停止、极限与指数计算}
\begin{exercise}
独立公平步长 $\xi_i\in\{-1,1\}$, $S_n=\sum_{i=1}^n\xi_i$, $S_0=0$. 求首次到达 $-a$ 或 $b$ 时到达下边界的概率及平均到达时间, 其中 $a,b$ 为正整数. 另用中心极限定理近似 $\Prob(S_{4\cdot10^6}>4000)$.
\end{exercise}
\begin{solution}
停止时间 $T=\inf\{n:S_n\in\{-a,b\}\}$ 几乎必然有限且有有限期望: 在每段 $a+b$ 步内, 全向右走的概率为 $2^{-(a+b)}$, 从任一内部位置出发都会先到边界. 所以 $\Prob(T>k(a+b))\le(1-2^{-(a+b)})^k$, 给出几何尾界.

$S_n$ 与 $S_n^2-n$ 都是鞅. 先对有界时间 $T\wedge n$ 使用可选停止, 得 $\E S_{T\wedge n}=0$ 和 $\E S_{T\wedge n}^2=\E(T\wedge n)$. 停止后位置始终在 $[-a,b]$, 可用有界控制令 $n\to\infty$; 时间部分则由单调收敛和上面的有限期望成立. 设 $p=\Prob(S_T=-a)$, 则 $-ap+b(1-p)=0$, 所以 $p=b/(a+b)$. 再代入第二式,
\[
\E T=a^2\frac b{a+b}+b^2\frac a{a+b}=ab.
\]
这里不能直接说``停止时间有限, 所以可选停止''. 真正支持极限的是停止位置有界及停时的尾估计.

固定大时间时, $S_n$ 均值0、方差 $n$. 本题标准差为2000, 阈值是两倍标准差, 所以正态近似为 $1-\Phi(2)\approx0.02275$. 中心极限定理给近似极限, 不给此有限 $n$ 的自动误差保证; 若要误差上界, 还需额外的定量定理.
\end{solution}

\begin{exercise}
$X_i$ 独立且可积, $\E X_i=0$. 在自然过滤下, 判断以下过程是否必为鞅:
$A_n=\sum_{i=1}^n iX_i$,
$B_n=\sum_{i=1}^nX_i^2-n$,
$C_n=\prod_{i=1}^n(1+X_i)$,
$D_n=\prod_{i=1}^n(X_i-1)$.
\end{exercise}
\begin{solution}
$A_n$ 的新增项 $(n+1)X_{n+1}$ 条件均值为0, 故是鞅. $B_n$ 不必可积, 即使可积, 新增项条件均值也为 $\E X_{n+1}^2-1$, 不一定为0. 取所有 $X_i=0$ 就有 $B_n=-n$, 给出反例. 若额外假定各二阶矩为1, 它才是鞅.

对 $C_n$, 有限个独立可积因子的绝对乘积期望等于绝对期望乘积, 所以可积; 条件期望为 $C_n\E(1+X_{n+1})=C_n$, 故是鞅. $D_n$ 的条件期望是 $-D_n$, 一般不同于自身; 同样取 $X_i=0$ 得 $D_n=(-1)^n$. 把它改成 $(-1)^nD_n=\prod_{i=1}^n(1-X_i)$ 则成为鞅.
\end{solution}

\begin{exercise}
$Z\sim\Normal(0,1)$, $\Phi$ 为标准正态分布函数. 对实数 $c,d$ 与 $a<b$, 求 $\E\ee^{cZ+d}$ 及 $\E[\ee^{cZ+d}\ind_{\{a<Z<b\}}]$.
\end{exercise}
\begin{solution}
配方是计算的核心: $cz-z^2/2=c^2/2-(z-c)^2/2$. 因而
\[
\begin{aligned}
\E[\ee^{cZ+d}\ind_{\{a<Z<b\}}]
&=\ee^{d+c^2/2}\int_a^b\frac1{\sqrt{2\pi}}\ee^{-(z-c)^2/2}\dd z\\
&=\ee^{d+c^2/2}\bigl[\Phi(b-c)-\Phi(a-c)\bigr].
\end{aligned}
\]
令区间扩大为整条实轴, 得 $\E\ee^{cZ+d}=\ee^{d+c^2/2}$. 用在对数正态股票价格时, 指数因子与阈值一起平移; 忘掉阈值平移就会把看涨期权的第一个正态概率算错.
\end{solution}

\originalsection{逐讲来源对应}
全部来源是官方18.600 Fall 2019讲义目录中的38份文件. 下表中的页数是原始PDF页数, 包括渐显重复、标题与来源页, 不是本中文稿页数. 第16、29讲无讲义. 前两次复习重复数学内容归入相应章节, 期末复习的训练目标由本章练习完整体现. 第39讲末的课程推荐与告别页没有新增数学结论, 不作为教材正文复制.

\begin{longtable}{r p{8.0cm} r}
\toprule
讲次 & 主要内容与中文对应 & 原PDF页数\\
\midrule
\endhead
1 & 排列组合; 第1章 & 81\\
2 & 多项式系数与计数; 第1章 & 53\\
3 & 样本空间与集合; 第1章 & 47\\
4 & 概率公理; 第1章 & 54\\
5 & 等可能模型; 第1章 & 42\\
6 & 条件概率; 第2章 & 39\\
7 & Bayes公式与独立; 第2章 & 62\\
8 & 离散随机变量; 第3章 & 53\\
9 & 离散期望; 第3章 & 62\\
10 & 方差; 第3章 & 73\\
11 & 二项分布与投资组合; 第3章 & 62\\
12 & 泊松分布; 第4章 & 65\\
13 & 泊松过程; 第4章 & 59\\
14 & 更多离散分布; 第4章 & 63\\
15 & 第一次期中复习; 第13章 & 139\\
17 & 连续随机变量; 第5章 & 103\\
18 & 正态分布; 第5章 & 56\\
19 & 指数分布; 第5章 & 70\\
20 & 更多连续分布; 第5章 & 59\\
21 & 联合分布; 第6章 & 69\\
22 & 独立变量之和; 第6章 & 44\\
23 & 和的期望与次序统计量; 第7章 & 61\\
24 & 协方差与条件练习; 第7章 & 70\\
25 & 条件期望; 第8章 & 59\\
26 & 矩母函数与特征函数; 第8章 & 75\\
27 & 弱大数定律; 第9章 & 52\\
28 & 第二次期中复习; 第13章 & 184\\
30 & 中心极限定理; 第9章 & 78\\
31 & 强大数定律与Jensen; 第9章 & 64\\
32 & 马尔可夫链; 第10章 & 54\\
33 & 熵; 第11章 & 58\\
34 & 鞅与可选停止; 第12章 & 72\\
35 & 鞅与风险中性概率; 第12章 & 80\\
36 & 风险中性概率与Black--Scholes; 第12章 & 62\\
37 & 期末复习I; 第13章 & 20\\
38 & 期末复习II; 第13章 & 15\\
39 & 期末复习III; 第13章 & 40\\
补充 & 鞅、风险中性概率与Black--Scholes; 第12章 & 14\\
\midrule
合计 & 37份主讲义与1份补充 & 2413\\
\bottomrule
\end{longtable}
逐文件原始直链、资源页、下载日期、版权核查与SHA-256见源码包的source-manifest.json. 同一主源只算一个具体学科, 不把旧课号18.440另行计数. 原课与本稿的关系是数学内容重构与补充, 不把本稿新增证明归为原作者已经写出的证明.
